%% notes-tex/essentiality-figure-gate/sections/3-models.tex
%% SOURCE: the macros of tables/gate_numbers.tex, written by
%% experiments/028-gene-essentiality/scripts/essentiality_gate_numbers.py from
%% results/fcl_comparison.json, family_ensembles_folds.json, dcell_style_folds.json and
%% genome_heldout_readout.json.

\clearpage
\section{The models and comparators on file\secstatus{todo}}
\label{sec:models}

\pfstep{Family 1, the transformer} The cell graph transformer over the nine-graph
perturbed gene graph with the pathway multi-hot and ProtT5 as node features, the deleted
gene's token pooled to a logit, and a zero-initialized affine that adds the metabolic
columns computed offline from yeast-GEM 9.0.2: the growth ratio of the deletion on one
medium, the deletion's parsimonious flux vector as $\log(1 + |v|)$, and six counts read
off the gene-protein-reaction rules. As a five-fold ensemble on the released split the
single-medium form reads \gfamoneMedium{} [\gfamoneMediumLo{}, \gfamoneMediumHi{}], the
rule-count form \gfamruleCounts{} [\gfamruleCountsLo{}, \gfamruleCountsHi{}], and the
layer-free solve alone \gfamlayerFree{} [\gfamlayerFreeLo{}, \gfamlayerFreeHi{}].

\pfstep{Family 2, graph convolution on GO} A two-layer graph convolution over the union
of nine gene-gene graphs on the propagated Gene Ontology multi-hot (2,362 terms over
5,501 annotated genes), which reads \gfamgoGcn{} [\gfamgoGcnLo{}, \gfamgoGcnHi{}] as a
five-fold ensemble. It is inductive in the regime key's sense: a gene's own annotation
and graph position, no call of any other gene.

\pfstep{Their average} Each family's logits standardized and weighted equally with nothing
fitted. The best of the \gnPairAverages{} two-family averages on file is
\gbestAverage{} [\gbestAverageLo{}, \gbestAverageHi{}] (\gbestAverageNames{}); the average
was fixed in advance before the \gnResplits{} repeated splits were read, where it beats
the layer-free form by \gresplitDiffaverageMinuslayerfree{}
(\gresplitDiffaverageMinuslayerfreeUp{} of \gresplitDiffaverageMinuslayerfreeOf{}) and
graph convolution on GO by \gresplitDiffaverageMinusGOconvolution{}
(\gresplitDiffaverageMinusGOconvolutionUp{} of \gresplitDiffaverageMinusGOconvolutionOf{}).

\pfstep{Comparators on the same genes} Flux Cone Learning (Merzbacher 2025) scores
\gfclAuroc{} on the \gnHeldOut{} genes; on its own released held-out set of
\gfclReleasedN{} it reads \gfclReleasedAuroc{}, and its published five-fold mean of
\gfclPublishedCv{} was recomputed at \gfclRecomputedCv{} from the released notebook, so the
comparison is against a reproduced baseline. FBA with yeast-GEM 9.0.2 and nothing fitted
scores \gfbaAuroc{} [\gfbaAurocLo{}, \gfbaAurocHi{}]. A DCell-style model over
\gdcellTerms{} GO terms, trained under the same fold protocol, reads \gdcellEnsemble{}
[\gdcellEnsembleLo{}, \gdcellEnsembleHi{}] as a fold ensemble (fold mean
\gdcellFoldMean{}).

\pfstep{The genome-wide limit} On the \ggenomeN{} genes with no GEM reaction the best
arm on file is \ggenBestName{} at \ggenBest{}; the transformer with the solve reads
\ggenlayerFree{}, graph convolution on ProtT5 over the union \ggengcnProt{}, graph
convolution on the physical graph's degree and clustering \ggengcnDegree{}, and the FBA
ratio \ggenfba{} (a constant). Every genome-wide arm is in \ref{sec:performance}.
